\chapter{Introduction}
\section{What this course will teach}
This workshop introduces you to Python coding, electronics, and project design. You will
build several projects, starting simple (fewer components and shorter code) and working up
to more complex ones (more components and more involved code).

Every project is built around a \textbf{microcontroller}: a small, complete computer on a
single chip. It has a processor, memory, and pins that can sense and control the outside
world. It is far less powerful than a laptop, but it starts instantly, uses very little
power, and can be built into almost anything: thermostats, game controllers, cars. The one
in your kit is an ESP32-C3, mounted on a small development board made by
\href{https://www.seeedstudio.com/Seeed-XIAO-ESP32C3-p-5431.html}{Seeed Studio}. It runs
\textbf{MicroPython}, a version of the Python language small enough to fit on a
microcontroller. Alongside it, you will use LEDs, buttons, a speaker, sensors, and a small
screen.

Use the written instructions and diagrams, but also feel free to experiment with the
components to see how they work and what they do.

\begin{tcolorbox}[colback=blue!5!white,colframe=blue!40!black]
    \textbf{New to Python or MicroPython?}
    If Python is new to you, or if you would like a quick refresher on how MicroPython interacts
    with hardware pins, check out \textbf{Appendix \ref{python_primer} (Python Primer)} before starting
    Project 1. If you already have Python experience, the primer's MicroPython section provides a
    handy overview of the hardware patterns used across all the projects.
\end{tcolorbox}

\section{Project Kit Contents}
\begin{figure}[H]
\centering
    \maybeincludegraphics[width=\linewidth]{common/case_and_parts.jpg}
    \caption{This image shows all of the contents of the project kit and labels what each one is}
\end{figure}

\section{Working with the microcontroller}
\subsection{Pinout Reference} \label{pinout}
A \textbf{pinout} is a map of what each pin on a chip or board does. Electronics engineers
keep one handy whenever they wire something up. Here is the pinout for the microcontroller
board used in this workshop:
\begin{figure}[H]
\centering
    \maybeincludegraphics[width=.6\linewidth]{common/esp32c3_pinout.png}
    \caption{A simplified pinout of the Seeed Studio ESP32-C3 development board}
\end{figure}

The label colors tell you what kind of pin each one is:
\begin{itemize}
    \item \textbf{Green (GPIO):} \textbf{General-Purpose Input/Output} pins. Your code decides what
    each one does: it can switch a pin on or off to control an LED or speaker, or read a pin to
    check a button or sensor.
    \item \textbf{Red (5V and 3V3):} Power pins that supply 5 volts or 3.3 volts to your circuits.
    \item \textbf{Black (GND):} Ground, the 0-volt return path that every circuit connects back to.
\end{itemize}

In the code throughout this guide, you will see pin numbers like \textbf{Pin(5)}. This means
\textbf{GPIO5} on the pinout, which is the 4th pin down on the left. If your code is not working
as expected, double check the pin numbers in the code against where your wires are plugged in.

\subsection{Project Guide Updates}

In case you are reading an out of date version of this guide, you can find the latest at
\url{https://netapp-ptc.github.io/YWIT-2026/project_guide.pdf}.

You can also find the source code for all of the workshops and other useful information in the GitHub
repository at \url{https://github.com/NetApp-PTC/YWIT-2026}.
